% preamble.tex (for mdframed example)
\usepackage{pdflscape}        % For landscape pages

% --- Fonts and Encoding ---

\usepackage{fontspec}         % Allows font specification
\usepackage{luatexja}          % For CJK scripts in LuaLaTeX, load after fontspec
% Set font for Chinese script
\newfontfamily\chinesefont{Noto Sans SC}
% Set font for Devanagari script
\newfontfamily\devanagarifont[Language=Sanskrit]{Noto Sans Devanagari}

% Main Font
\setmainfont{Times New Roman}

% Math Font
\setmathfont{XITS Math}

% Monospace Font
\setmonofont{Courier New}

% --- Packages for Tables ---

\usepackage{longtable}         % For long tables
\usepackage{array}             % For table column width specification
\usepackage{booktabs}         % For table rules
\usepackage{ragged2e}         % For text alignment (used with \newcolumntype and footnotes)

% --- Other Packages ---

\usepackage{mdframed}
\usepackage[x11names]{xcolor}    % Required for specifying custom colors, load before tcolorbox
\usepackage[version=4]{mhchem}  % Formula subscripts using \ce{}

% --- Color Definitions ---

\definecolor{headingblue}{RGB}{23,48,191}
\definecolor{boxtitle}{HTML}{F0F4F8}
\definecolor{boxbody}{HTML}{FBFDFF}
\definecolor{mainboxframe}{HTML}{F0F4F8}
\definecolor{subboxframe}{HTML}{F0F4F8}

% --- mdframed Styles ---

\newmdenv[
  skipabove=12pt,
  skipbelow=12pt,
  leftmargin=0pt,
  rightmargin=0pt,
  innertopmargin=2mm,
  innerbottommargin=2mm,
  innerleftmargin=2mm,
  innerrightmargin=2mm,
  backgroundcolor=boxbody,
  frametitlebackgroundcolor=boxtitle,
  frametitlefont={\bfseries\sffamily},
  frametitleaboveskip=\dimexpr-\ht\strutbox+2pt\relax,
  frametitlealignment=\raggedright,
  roundcorner=0pt,
  nobreak=false,
  usetwoside=false,
  splittopskip=2\baselineskip,
  splitbottomskip=\baselineskip,
  linewidth=0pt,
  frametitle={#1}
]{mainbox}

\newmdenv[
  skipabove=12pt,
  skipbelow=12pt,
  leftmargin=0pt,
  rightmargin=0pt,
  innertopmargin=2mm,
  innerbottommargin=2mm,
  innerleftmargin=2mm,
  innerrightmargin=2mm,
  backgroundcolor=boxbody,
  frametitlebackgroundcolor=boxtitle,
  frametitlefont={\bfseries\sffamily},
  frametitleaboveskip=\dimexpr-\ht\strutbox+2pt\relax,
  frametitlealignment=\raggedright,
  roundcorner=0pt,
  nobreak=false,
  usetwoside=false,
  splittopskip=2\baselineskip,
  splitbottomskip=\baselineskip,
  linewidth=0pt,
  frametitle={#1}
]{subbox}

% --- Custom Column Type (using ragged2e) ---

\newcolumntype{R}[1]{>{\RaggedRight}p{#1}}

% --- Epigraph ---

\usepackage{epigraph}
\setlength\epigraphwidth{.9\textwidth}
\newenvironment{quotepara}
  {\setlength{\parindent}{1em}\setlength{\parskip}{0.6ex}\itshape\raggedright}
  {}
\renewcommand{\textflush}{quotepara}

% --- Small Caps ---

\newcommand{\flatcaps}[1]{\textsc{\MakeLowercase{#1}}}

% --- Custom Chapter/Section Styles ---

\usepackage[compact]{titlesec}  % Allows creating custom chapter styles
\titleformat{\chapter}[display]
  {\fontsize{60}{62}\bfseries}
  {\thechapter}
  {0pt}
  {\huge\noindent}
\titlespacing*{\chapter}{0pt}{0pt}{40pt}

\titleformat{\section}[block]
  {\large\bfseries}
  {\thesection}
  {1em}
  {}
\titlespacing*{\section}{0pt}{2\baselineskip}{1\baselineskip}

\titleformat{\subsection}
  {\large\bfseries}
  {\thesubsection}
  {0.6em}
  {\flatcaps}
\titlespacing*{\subsection}{0pt}{1\baselineskip}{1\baselineskip}

\titleformat{\subsubsection}[runin]
  {}
  {\thesubsubsection}
  {0.6em}
  {}
  {\itshape}
\titlespacing*{\subsubsection}{0pt}{8pt}{5pt}

\titleformat{\paragraph}[runin]
  {\flatcaps}
  {\theparagraph}
  {0pt}
  {}

% --- Footnotes ---

\usepackage[ragged,hang]{footmisc}  % Load footmisc with ragged option
\renewcommand{\footnotelayout}{\RaggedRight\footnotesize} % Typeset footnotes in \RaggedRight
\setlength{\footnotemargin}{1.5em}   % Adjust space between number and text
\makeatletter
\renewcommand{\@makefntext}[1]{%
    \parindent 1em%                    Set parindent for footnote text
    \noindent
    \hb@xt@ 1.8em{%                   Set hanging indent for footnote text
        \hss\@thefnmark.%
    }
    \RaggedRight #1%                  Typeset footnote text ragged right
}
\makeatother

% --- Captions ---

\usepackage{caption}
\captionsetup{
  font={small},
  labelfont={bf},
  textfont={},
  width=0.9\textwidth,
  justification=justified,
  labelformat=default,
  labelsep=period,
  format=plain
}
\renewcommand{\captionlabelfont}{\bfseries\scshape}

% --- Hyperlinks ---

\usepackage{hyperref}           % Load after most other packages, but before cleveref
\hypersetup{
  colorlinks = true,
  linkcolor = {headingblue},
  citecolor = {headingblue},
  urlcolor = {headingblue}
}

% --- Miscellaneous ---

\usepackage{iftex}             % Detects the engine used
\usepackage{microtype}         % Improves typography
\usepackage{etoolbox}
%% Use lining fonts
\newcommand\lining{\addfontfeatures{Numbers={Monospaced, Lining}}}
\AtBeginEnvironment{tabular}{\lining} % In tables
\renewcommand{\theequation}{ {\lining\arabic{equation}}} % For equation numbers

% --- Index (if needed) ---

\usepackage{makeidx}
\makeindex

% --- Other Typography Settings ---

\frenchspacing                  % Single space after periods

% --- Microtype Settings (adjust only if needed) ---

% \microtypesetup{
%    tracking = true,
%    factor = 1100,
%    stretch = 15,
%    shrink = 15
% }
